\documentclass[12pt]{article}
\usepackage[utf8]{inputenc}
\usepackage[T1]{fontenc}
\usepackage{hyperref}
\title{Tai Yo Language Resources for K-12 Curriculum Development}
\author{Compiled by Claude Sonnet 5 (claude-sonnet-5) for Andr\'as Kornai}
\date{}
\begin{document}
\maketitle

\section{Introduction}
Tai Yo (ISO 639-3 tyj; also Tai Mène, Nyaw; in Vietnam Tai Do, Tai Quy Chau) is a Northern Tai language of Vietnam and Laos. It is closely related to Tai Pao. In Laos it is spoken in Khamkeut District, Borikhamxay Province. Wikipedia cites about 60,000 speakers from a 1990--1995 census, while the task description gives about 300, so the figures are inconsistent. A distinctive Tai Yo script existed but is no longer in use. The search found almost nothing beyond encyclopedic and academic pages.

\section{Primary Educational Resources}
None found.

\section{Linguistic and Reference Materials}
\begin{itemize}
\item Wikipedia article on Tai Yo: \url{https://en.wikipedia.org/wiki/Tai_Yo_language}
\item Wikipedia article on the related Tai Pao language: \url{https://en.wikipedia.org/wiki/Tai_Pao_language}
\item Papers by James R. Chamberlain, who has published on Tai languages of the region; I did not verify which papers concern Tai Yo: \url{https://nuol.academia.edu/JamesChamberlain}
\end{itemize}

\section{Religious and Community Resources}
None found.

\section{Digital and Audio Resources}
None found.

\section{Curriculum Development Context}
There is no evidence of Tai Yo being taught or of any teaching material. Related Lao-language education in Laos exists, but that is a different language. Any curriculum would start from documentation work with the community, including possible revival of the historical script.
\end{document}